\documentclass[11pt]{article}

\usepackage[margin=1in]{geometry}
\usepackage{amsmath, amssymb}
\usepackage{xcolor}
\usepackage{listings}
\usepackage{hyperref}
\usepackage{array}
\usepackage{booktabs}

\hypersetup{
    colorlinks=true,
    linkcolor=blue!60!black,
    urlcolor=blue!60!black,
    citecolor=blue!60!black
}

\lstdefinestyle{bugpython}{
  language=Python,
  basicstyle=\ttfamily\small,
  keywordstyle=\color{blue!60!black},
  stringstyle=\color{green!40!black},
  commentstyle=\color{gray!70!black},
  showstringspaces=false,
  breaklines=true,
  frame=single,
  rulecolor=\color{gray!30},
  columns=fullflexible,
  keepspaces=true
}

\title{SymPy Bug Report}
\author{}
\date{}

\begin{document}

\maketitle

\section{Bug 47: linsolve misses the a = 0 branch}

\subsection{Minimal Reproducer}

\medskip

\begin{lstlisting}[style=bugpython]
from sympy import S, linsolve, symbols

x, a = symbols("x a")
sol = linsolve([a*x - a], [x])
print("solution:", sol)
print("equation at a=0, x=2:", (a*x - a).subs({a: 0, x: 2}))
print("reported contains x=2 after a=0:", (S(2),) in sol.subs(a, 0))
\end{lstlisting}

\subsection{Observed Behavior}

\medskip

\begin{lstlisting}[style=bugpython]
SymPy version: 1.14.0
solution: {(1,)}
equation at a=0, x=2: 0
reported contains x=2 after a=0: False
\end{lstlisting}

SymPy returns only the singleton solution \(\{(1,)\}\).  However, after
specializing the parameter to \(a=0\), the original equation is satisfied by
\(x=2\), while \((2,)\) is not contained in the specialized reported solution.
The returned set therefore omits valid solutions on the degenerate parameter
branch.

\subsection{Expected Behavior}

The mathematically correct answer is conditional: \(x=1\) when \(a\ne 0\), and
all values of \(x\) when \(a=0\).  Equivalently, the solution set should include
the generic branch
\[
  a \ne 0,\qquad x=1,
\]
and the degenerate branch
\[
  a=0,\qquad x \text{ arbitrary}.
\]

\subsection{Mathematical Explanation}

The equation supplied to \texttt{linsolve} is
\[
  a x - a = 0,
\]
or, after factoring,
\[
  a(x-1)=0.
\]
If \(a\ne 0\), division by \(a\) is valid and the equation reduces to
\(x-1=0\), so \(x=1\).  If \(a=0\), the equation instead becomes
\[
  0\cdot (x-1)=0,
\]
which is the identity \(0=0\) and imposes no constraint on \(x\).  Thus values
such as \(x=2\) are valid solutions on the \(a=0\) branch.

The result \(\{(1,)\}\) is not an equivalent representation of the full
solution set.  Substituting \(a=0\) into the original equation and \(x=2\) gives
zero, but substituting \(a=0\) into the reported solution leaves only
\(\{(1,)\}\), which excludes \((2,)\).  The missing branch is therefore a real
loss of solutions, not a harmless simplification or alternate set notation.

\subsection{Source-Level Diagnosis}

The diagnosis identifies the sparse linear-solver domain and row-reduction
pipeline as the source-level mechanism.  The public call
\texttt{linsolve([a*x - a], [x])} is handled in
\nolinkurl{sympy/solvers/solveset.py}, where \texttt{linsolve} delegates list
inputs to \texttt{\_linsolve}.  That internal solver is imported from
\nolinkurl{sympy/polys/matrices/linsolve.py}.  In the traced path,
\texttt{\_linsolve} converts the equation to coefficient dictionaries, builds a
sparse augmented matrix, and computes its reduced row echelon form with
\texttt{sdm\_irref}.

For the representative equation, the linear conversion gives coefficient
\(\{x: a\}\) and constant \(-a\).  The matrix-domain construction in
\nolinkurl{sympy/polys/matrices/linsolve.py} uses the elements of this
augmented row to construct the coefficient domain as the rational-function
field \(\mathbb{Z}(a)\).  In that field, \(a\) is treated as a nonzero
invertible coefficient unless it is the zero rational function.  Row reduction
therefore divides the row \([a,a]\) by \(a\), obtains the normalized row
\([1,1]\), and returns the generic solution \(x=1\).

\begin{lstlisting}[style=bugpython]
K, elems_K = construct_domain(elems, field=True, extension=True)
...
Arref, pivots, nzcols = sdm_irref(Aaug)
\end{lstlisting}

This source-level behavior explains the mathematical failure: the solver has
effectively solved the generic branch where \(a\ne 0\), but the returned
\texttt{FiniteSet} carries no condition \(a\ne 0\) and does not split off the
specialization \(a=0\).  The diagnosis presents a plausible fix direction:
for systems with symbolic coefficients, \texttt{linsolve} should either
document and return only generic solutions, or carry pivot nonzero conditions
and split degenerate parameter branches.  In this case, branch-aware handling
would return \(x=1\) for \(a\ne 0\) and all \(x\) for \(a=0\).

\subsection{Additional Failing Instances}

The independent verification checks the family
\[
  a(x-r)=0
\]
for integer roots \(r=1,2,3,4,5,6\).  In each case, the generic solution is
\(x=r\), but the specialization \(a=0\) turns the equation into \(0=0\), so
every \(x\) is a solution.  The verification chooses \(x=r+1\) as an explicit
omitted solution on the \(a=0\) branch.

\begin{center}
\renewcommand{\arraystretch}{1.3}
\begin{tabular}{@{}>{\raggedright\arraybackslash}p{0.44\linewidth}>{\raggedright\arraybackslash}p{0.23\linewidth}>{\raggedright\arraybackslash}p{0.23\linewidth}@{}}
\toprule
\textbf{Input / Expression} & \textbf{SymPy behavior} & \textbf{Expected behavior} \\
\midrule
\(\operatorname{linsolve}([a(x-1)],[x])\) & Returns \(\{(1,)\}\); after \(a=0\), omits \((2,)\) & \(x=1\) if \(a\ne 0\); all \(x\) if \(a=0\) \\
\(\operatorname{linsolve}([a(x-2)],[x])\) & Returns \(\{(2,)\}\); after \(a=0\), omits \((3,)\) & \(x=2\) if \(a\ne 0\); all \(x\) if \(a=0\) \\
\(\operatorname{linsolve}([a(x-3)],[x])\) & Returns \(\{(3,)\}\); after \(a=0\), omits \((4,)\) & \(x=3\) if \(a\ne 0\); all \(x\) if \(a=0\) \\
\(\operatorname{linsolve}([a(x-4)],[x])\) & Returns \(\{(4,)\}\); after \(a=0\), omits \((5,)\) & \(x=4\) if \(a\ne 0\); all \(x\) if \(a=0\) \\
\(\operatorname{linsolve}([a(x-5)],[x])\) & Returns \(\{(5,)\}\); after \(a=0\), omits \((6,)\) & \(x=5\) if \(a\ne 0\); all \(x\) if \(a=0\) \\
\(\operatorname{linsolve}([a(x-6)],[x])\) & Returns \(\{(6,)\}\); after \(a=0\), omits \((7,)\) & \(x=6\) if \(a\ne 0\); all \(x\) if \(a=0\) \\
\bottomrule
\end{tabular}
\end{center}

The same reasoning applies uniformly across the family.  The factor \(a\) is
the only coefficient carrying the parameter branch, so specializing \(a=0\)
annihilates the equation regardless of the integer root \(r\).

\subsection{Related Issues Assessment}

The bug-finding harness performed a best-effort deduplication check.  This
specific failure mode is documented in the open issue
\href{https://github.com/sympy/sympy/issues/16861}{sympy/sympy\#16861},
``Proper solution of linear equations with symbolic coefficients'' (open, 18
comments, opened 2019).  That issue raises the same symbolic-coefficient
missing-\(a=0\)-branch problem using the essentially identical example
\(a x = b\) (this bug's \(a x - a\) is the \(b=a\) case), observes that the
solve/solveset/linsolve solutions ``are not valid for a=0'', and proposes a
conditional \texttt{linsolve} variant that carries the degenerate \(a=0\)
branch.  Consistent with this, the SymPy documentation shows
\texttt{linsolve} solving systems with symbolic coefficients by producing
generic determinant-denominator formulas without splitting singular-parameter
cases, which is its long-standing documented generic-coefficient convention.
The deduplication verdict is: family known, specific instance.  The exact
reproducer \texttt{linsolve([a*x - a], [x])} and the result \(\{(1,)\}\) fall
squarely within the failure mode \#16861 describes and remain individually
actionable as a regression test, but this is a known, long-standing limitation
rather than a previously unreported failure mode.

\subsection{Proposed SymPy Regression Test}

\medskip

\begin{lstlisting}[style=bugpython]
import pytest
from sympy import S, linsolve, symbols

x, a = symbols("x a")

@pytest.mark.parametrize("root", [1, 2, 3, 4, 5, 6])
def test_linsolve_keeps_zero_parameter_branch(root):
    expr = a*(x - root)
    sol = linsolve([expr], [x])
    extra_solution = S(root + 1)
    assert (extra_solution,) in sol.subs(a, 0)
\end{lstlisting}

\end{document}
